\documentclass[11pt]{article}
\usepackage[T1]{fontenc}
\usepackage{lmodern}
\usepackage[a4paper,margin=28mm]{geometry}
\usepackage{amsmath,amssymb,amsthm,mathtools}
\usepackage{microtype}
\usepackage[hidelinks]{hyperref}
\usepackage{enumitem}
\newtheorem{theorem}{Theorem}[section]
\newtheorem{lemma}[theorem]{Lemma}
\newtheorem{proposition}[theorem]{Proposition}
\newtheorem{corollary}[theorem]{Corollary}
\theoremstyle{remark}
\newtheorem{remark}[theorem]{Remark}
\newcommand{\Prob}{\mathcal P}
\newcommand{\E}{\mathbb E}
\newcommand{\R}{\mathbb R}
\newcommand{\dd}{\,d}
\newcommand{\push}{_{\#}}
\DeclareMathOperator{\supp}{supp}
\title{The interval minimizes negative-exponent\protect\\ Riesz capacity ratios}
\author{GPT-6 Astra}
\date{November 2, 2026}
\begin{document}
\maketitle

\begin{abstract}
For a compact set $K\subset\R$ containing at least two points,
let $C_t(K)$ be the $t$th root of the
maximum of $\iint |x-y|^t\dd\mu(x)\dd\mu(y)$ over probability measures on $K$.
We prove that, for every $0<r<s<1$, the ratio $C_s(K)/C_r(K)$ is minimized
by an interval, with equality only for nondegenerate closed intervals.
This resolves the interval conjecture of Clark and Laugesen, restated by Fan.
The main step is a comparison of energies of monotone images of beta measures.
After a logistic change of coordinates, the relevant pair distribution is a
positive mixture of cells from Poisson partitions. A likelihood-ratio
comparison orders the mixing rates, and refinement of the partitions gives
the energy comparison by concavity. The argument also derives the beta
equilibrium measure and the interval capacity formula.
\end{abstract}

\section{The capacity inequality}

For a Borel probability measure $\mu$ with compact support, write
\[
 \mathcal E_t(\mu)=\iint |x-y|^t\dd\mu(x)\dd\mu(y),\qquad t>0.
\]
For a nonempty compact set $K\subset\R$, define
\[
 U_t(K)=\max_{\mu\in\Prob(K)}\mathcal E_t(\mu),
 \qquad C_t(K)=U_t(K)^{1/t}.
\]
The maximum exists: the probability measures on $K$ form a weakly compact
space, and the bounded continuous kernel makes the energy continuous.
For the latter assertion, weak convergence of the product measures follows
by uniformly approximating continuous functions on $K\times K$ by finite
sums of functions $f(x)g(y)$. If $K$ contains two distinct points, placing
half the mass at each shows that $U_t(K)>0$.

\begin{theorem}\label{thm:main}
Let $K\subset\R$ be compact and contain at least two points. For every
$0<r<s<1$,
\begin{equation}\label{eq:main}
 \frac{C_s(K)}{C_r(K)}\geq
 \frac{C_s([-1,1])}{C_r([-1,1])},
 \qquad
 C_t([-1,1])=
 \left[\frac{\Gamma((1+t)/2)\Gamma(1-t/2)}{\sqrt\pi}\right]^{1/t}
 \quad(0<t<1).
\end{equation}
Equality holds if and only if $K$ is a nondegenerate closed interval.
\end{theorem}

In the conventional Riesz notation, $C_t=\operatorname{Cap}_{-t}$.
Thus the inequality is the one-dimensional negative-exponent case
$p=-s$, $q=-r$ of the interval conjecture in Clark and
Laugesen~\cite[Conjecture 12]{CL}; Fan states exactly this formulation
as Conjecture 4.1~\cite{Fan}. Fan's Theorem 4.3 proves positivity of
equilibrium masses for finite subsets of the line, supporting numerical
tests of the conjecture. Our argument applies directly to every compact
set, including infinite and disconnected sets.

We first prove a comparison for monotone maps. For $0<t<1$, set
\begin{equation}\label{eq:beta}
 a_t=\frac{1-t}{2},\qquad
 d\beta_t(x)=\frac{[x(1-x)]^{a_t-1}}{B(a_t,a_t)}\dd x,
 \quad 0<x<1,
 \qquad e_t=\mathcal E_t(\beta_t).
\end{equation}
Here $B(a,b)=\Gamma(a)\Gamma(b)/\Gamma(a+b)$. In particular $e_t>0$.
We will derive, rather than assume, that $e_t=U_t([0,1])$.
For a bounded nondecreasing map $T:(0,1)\to\R$, define
\begin{equation}\label{eq:normalized}
 H_t(T)=\left[\frac{\mathcal E_t(T\push\beta_t)}{e_t}\right]^{1/t}.
\end{equation}
The pushforward $T\push\beta_t$ is the distribution of $T(X)$ when
$X$ has distribution $\beta_t$.

\begin{theorem}[Monotone transport comparison]\label{thm:transport}
If $T:(0,1)\to\R$ is bounded and nondecreasing and $0<r<s<1$, then
\begin{equation}\label{eq:transport}
 H_r(T)\leq H_s(T).
\end{equation}
Equality holds if and only if $T(x)=A+Bx$ for Lebesgue-almost every $x$,
with $B\geq0$.
\end{theorem}

\section{Random intervals and averaging}

For $0<x<y<1$, let
\[
 D_T(x,y)=\frac{T(y)-T(x)}{y-x}.
\]
The ordered pair distribution
\begin{equation}\label{eq:pair}
 dP_t(x,y)=\frac{2}{e_t}(y-x)^t\dd\beta_t(x)\dd\beta_t(y),
 \qquad 0<x<y<1,
\end{equation}
is a probability measure. The factor 2 accounts for both orders of a
pair of independent beta variables. Directly from the definitions,
\begin{equation}\label{eq:moment}
 H_t(T)^t=\int D_T(x,y)^t\dd P_t(x,y).
\end{equation}

Let
\begin{equation}\label{eq:logistic}
 h(u)=\frac{1+\tanh u}{2},\qquad
 \Delta(u,z)=h(u+z)-h(u),\qquad u\in\R,\ z>0.
\end{equation}
We use Poisson partitions only through the elementary properties proved
below. A homogeneous Poisson process of rate $\lambda>0$ on $\R$ has
independent counts on disjoint intervals, and the count in an interval
of length $L$ has Poisson distribution with mean $\lambda L$.
Its consecutive points partition $\R$ into bounded cells almost surely.
Applying $h$ gives a countable partition $\mathcal Q_\lambda$ of $(0,1)$,
apart from its cell endpoints.

\begin{lemma}[The cell containing a uniform point]\label{lem:cell}
Choose $V$ uniformly in $(0,1)$, independently of $\mathcal Q_\lambda$.
Write the cell containing $V$ as $(h(u),h(u+z))$. The law of $(u,z)$ is
\begin{equation}\label{eq:cell}
 \lambda^2 e^{-\lambda z}\Delta(u,z)\dd u\dd z.
\end{equation}
\end{lemma}

\begin{proof}
The variable $v=h^{-1}(V)$ has density $h'(v)$ on $\R$ and is independent
of the process. Conditional on $v$, the distances to the nearest Poisson
points on its left and right are independent exponential variables of
rate $\lambda$: each survival probability is the probability of no
Poisson point in the corresponding interval. Thus, for endpoints $u$
and $u+z$ with $u<v<u+z$, their joint density is
$\lambda^2e^{-\lambda z}\dd u\dd z$.
Integrating over the sampled point gives
\[
 \lambda^2e^{-\lambda z}\int_u^{u+z}h'(v)\dd v
 =\lambda^2e^{-\lambda z}\Delta(u,z),
\]
which proves~\eqref{eq:cell}.

For an explicit normalization check, $h'\geq0$ and $\int_\R h'=1$, so
Tonelli's theorem gives
\begin{equation}\label{eq:gap-integral}
 \int_\R\Delta(u,z)\dd u
 =\int_\R\int_0^z h'(u+w)\dd w\dd u=z.
\end{equation}
Consequently the integral of~\eqref{eq:cell} is
$\int_0^\infty\lambda^2z e^{-\lambda z}\dd z=1$.
\end{proof}

For a bounded nondecreasing $T$, its one-sided limits exist, and
\[
 L_T=T(1-)-T(0+)\geq0.
\]
Changing $T$ at its at most countably many discontinuities has no effect
on any energy in this paper. Take a right-continuous version, and let
$\nu$ be its finite Lebesgue--Stieltjes measure on $(0,1)$.
Then $\nu((0,1))=L_T$.
Almost surely no Poisson endpoint is an atom of $\nu$, since $\nu$ has
at most countably many atoms and each fixed location has probability
zero of being a Poisson point.
For a cell $J=(x,y)$ with such endpoints, put
\[
 S_J=\frac{\nu(J)}{|J|}=D_T(x,y),
 \qquad A_{\mathcal Q}(v)=S_J\quad(v\in J).
\]
For almost every partition, countable additivity gives
\begin{equation}\label{eq:first-moment}
 \int_0^1 A_{\mathcal Q}(v)\dd v
 =\sum_{J\in\mathcal Q}\nu(J)=L_T.
\end{equation}
In particular, for $0<p<1$, Jensen's inequality yields
\begin{equation}\label{eq:integrability}
 0\leq\int_0^1 A_{\mathcal Q}(v)^p\dd v\leq L_T^p.
\end{equation}
This also covers singular Stieltjes measures and maps with jumps.

\begin{lemma}[Refinement comparison]\label{lem:refinement}
If $0<\lambda_1\leq\lambda_2$, the Poisson partitions can be coupled
so that $\mathcal Q_{\lambda_2}$ refines $\mathcal Q_{\lambda_1}$.
For every $0<p<1$ and bounded nondecreasing $T$,
\begin{equation}\label{eq:refinement}
 \E\int_0^1 A_{\mathcal Q_{\lambda_1}}(v)^p\dd v
 \geq
 \E\int_0^1 A_{\mathcal Q_{\lambda_2}}(v)^p\dd v.
\end{equation}
\end{lemma}

\begin{proof}
Superpose a rate-$\lambda_1$ process and an independent
rate-$(\lambda_2-\lambda_1)$ process, omitting the second process when
the rates are equal. Their union is a rate-$\lambda_2$ process by the
addition rule for independent Poisson counts.
Every coarse cell is bounded in the $u$ coordinate and hence contains
only finitely many added points almost surely.
If it splits into fine cells $J_1,\ldots,J_m$, its slope is
\[
 S_J=\sum_{i=1}^m\frac{|J_i|}{|J|}S_{J_i}.
\]
Concavity of $z^p$ gives
$|J|S_J^p\geq\sum_i|J_i|S_{J_i}^p$.
Summing over the countably many coarse cells, then taking expectations,
proves~\eqref{eq:refinement}. The sums are nonnegative and bounded
by~\eqref{eq:integrability}.
\end{proof}

\section{The beta pair distribution as a Poisson mixture}

\begin{lemma}[Mixture representation]\label{lem:mixture}
Under the change of variables $x=h(u)$, $y=h(u+z)$, the probability
measure $P_t$ in~\eqref{eq:pair} has density
\begin{equation}\label{eq:logistic-pair}
 \frac{\Delta(u,z)(\sinh z)^{t-1}}{Z_t}\dd u\dd z,
 \qquad
 Z_t=\int_0^\infty z(\sinh z)^{t-1}\dd z.
\end{equation}
Let $\alpha=1-t$ and
\begin{equation}\label{eq:weights}
 b_n(\alpha)=\frac{(\alpha)_n}{n!},\quad
 \lambda_n(\alpha)=2n+\alpha,\quad
 W(\alpha)=\sum_{n=0}^\infty\frac{b_n(\alpha)}{\lambda_n(\alpha)^2},
 \quad
 w_n(\alpha)=\frac{b_n(\alpha)}{\lambda_n(\alpha)^2W(\alpha)}.
\end{equation}
Here $(\alpha)_0=1$ and $(\alpha)_n=\alpha(\alpha+1)\cdots(\alpha+n-1)$.
Then $P_t$ is the mixture of the cell laws~\eqref{eq:cell}, with rate
$\lambda_n(\alpha)$ chosen with probability $w_n(\alpha)$.
\end{lemma}

\begin{proof}
Use
\[
 h'(u)=\frac{1}{2\cosh^2u},\qquad
 h(u)(1-h(u))=\frac{1}{4\cosh^2u},\qquad
 \Delta(u,z)=\frac{\sinh z}{2\cosh u\cosh(u+z)}.
\]
The unnormalized ordered-pair density transforms according to
\begin{align*}
 &(y-x)^t[x(1-x)y(1-y)]^{-(1+t)/2}\dd x\dd y\\
 &\hspace{18mm}=2^{t+1}(\sinh z)^{t-1}\Delta(u,z)\dd u\dd z.
\end{align*}
Equation~\eqref{eq:gap-integral} supplies the normalization in
\eqref{eq:logistic-pair}. The integral $Z_t$ is finite: its integrand
is of order $z^t$ at zero and decays exponentially times $z$ at infinity.

For every $z>0$, the binomial series gives the positive expansion
\begin{equation}\label{eq:binomial}
 (\sinh z)^{-\alpha}
 =2^\alpha e^{-\alpha z}(1-e^{-2z})^{-\alpha}
 =2^\alpha\sum_{n=0}^\infty b_n(\alpha)e^{-\lambda_n(\alpha)z}.
\end{equation}
Tonelli's theorem and~\eqref{eq:gap-integral} imply
$Z_t=2^\alpha W(\alpha)$, in particular $0<W(\alpha)<\infty$.
Substitution into~\eqref{eq:logistic-pair} expresses that density as
\[
 \sum_{n=0}^\infty w_n(\alpha)\,
 \lambda_n(\alpha)^2e^{-\lambda_n(\alpha)z}\Delta(u,z).
\]
Lemma~\ref{lem:cell} identifies each summand's normalized density.
\end{proof}

\begin{lemma}[Ordering the mixing rates]\label{lem:rates}
If $0<\beta<\alpha<1$, random integers with distributions
$w_n(\beta)$ and $w_n(\alpha)$ can be coupled as
$N_\beta\leq N_\alpha$. In this coupling,
\[
 2N_\beta+\beta<2N_\alpha+\alpha.
\]
\end{lemma}

\begin{proof}
Apart from a positive factor independent of $n$, the likelihood ratio is
\begin{equation}\label{eq:likelihood}
 \frac{w_n(\beta)}{w_n(\alpha)}
 =\frac{W(\alpha)}{W(\beta)}
   \frac{(\beta)_n}{(\alpha)_n}
   \left(\frac{2n+\alpha}{2n+\beta}\right)^2.
\end{equation}
The Pochhammer ratio decreases with $n$, since the ratio of successive
terms is $(n+\beta)/(n+\alpha)<1$. The last factor also decreases,
as $(2n+\alpha)/(2n+\beta)=1+(\alpha-\beta)/(2n+\beta)$.
Thus the likelihood ratio decreases.

For completeness, put $q_n=w_n(\beta)$ and $p_n=w_n(\alpha)$.
For each integer $k\geq0$, $q_i p_j\geq p_i q_j$ whenever
$i\leq k<j$. Summing shows
\[
 \left(\sum_{i\leq k}q_i\right)\left(1-\sum_{i\leq k}p_i\right)
 \geq
 \left(\sum_{i\leq k}p_i\right)\left(1-\sum_{i\leq k}q_i\right),
\]
and hence $\sum_{i\leq k}q_i\geq\sum_{i\leq k}p_i$.
Applying the two discrete inverse distribution functions to the same
uniform random number gives $N_\beta\leq N_\alpha$.
The strict rate inequality follows from $\beta<\alpha$.
\end{proof}

\begin{proposition}[Concave moments of chord slopes]\label{prop:moments}
For $0<r<s<1$, every bounded nondecreasing $T$, and every $0<p<1$,
\begin{equation}\label{eq:concave-moments}
 \int D_T^p\dd P_s\geq\int D_T^p\dd P_r.
\end{equation}
Also, for every $0<t<1$,
\begin{equation}\label{eq:mean-slope}
 \int D_T\dd P_t=L_T,\qquad
 \int D_T^p\dd P_t\leq L_T^p\quad(0<p<1).
\end{equation}
\end{proposition}

\begin{proof}
Take $\alpha=1-r$ and $\beta=1-s$ in Lemma~\ref{lem:rates}.
Conditional on the coupled rates, use Lemma~\ref{lem:refinement} to
couple the smaller-rate partition as a coarsening of the larger-rate
partition. Choose the same independent uniform point $V$ to select
the cells. Lemma~\ref{lem:mixture} identifies their endpoint laws as
$P_s$ and $P_r$, respectively. Conditional Jensen, as in
\eqref{eq:refinement}, proves~\eqref{eq:concave-moments} after averaging
over the rates. Equations~\eqref{eq:first-moment} and
\eqref{eq:integrability}, averaged over the mixture, prove
\eqref{eq:mean-slope}.
\end{proof}

\begin{proof}[Proof of Theorem~\ref{thm:transport}]
By the power-mean inequality on the probability space $P_s$, followed
by Proposition~\ref{prop:moments} with $p=r$,
\begin{equation}\label{eq:chain}
 \left(\int D_T^s\dd P_s\right)^{1/s}
 \geq\left(\int D_T^r\dd P_s\right)^{1/r}
 \geq\left(\int D_T^r\dd P_r\right)^{1/r}.
\end{equation}
All moments are finite by~\eqref{eq:mean-slope}.
Together with~\eqref{eq:moment}, this proves the inequality.

If equality holds, equality in the first step of~\eqref{eq:chain}
implies that $D_T$ is constant $P_s$-almost surely; this follows also
from strict concavity of $z^{r/s}$ applied to $D_T^s$.
The density of $P_s$ is strictly positive on $0<x<y<1$.
It follows that $T(y)-T(x)=B(y-x)$ for Lebesgue-almost every ordered
pair, for a constant $B\geq0$. By symmetry this holds for almost every
pair in $(0,1)^2$. Fubini's theorem then implies that $T(x)-Bx$ is
almost everywhere constant. Conversely, if $T(x)=A+Bx$ almost
everywhere, then $D_T=B$ almost everywhere, and $H_t(T)=B$ for every
$t$. This includes constant maps.
\end{proof}

\section{Equilibrium of the interval and arbitrary compact sets}

\begin{lemma}[Monotone realization of a measure]\label{lem:quantile}
Let $\mu$ be a probability measure supported on a compact subset of
$\R$, and let $0<t<1$. There is a bounded nondecreasing Borel map
$T:(0,1)\to\R$ such that $T\push\beta_t=\mu$.
For every $0<u<1$, $T\push\beta_u$ is supported on $\supp\mu$.
\end{lemma}

\begin{proof}
Let $F_t$ be the continuous strictly increasing distribution function
of $\beta_t$. It maps a beta-distributed variable to a uniform variable.
If $Q_\mu(v)=\inf\{x:\mu(( -\infty,x])\geq v\}$ is the quantile
function, take $T=Q_\mu\circ F_t$. The usual inverse-distribution
identity gives $T\push\beta_t=\mu$. In particular,
$T(x)\in\supp\mu$ for $\beta_t$-almost every $x$.
All the measures $\beta_u$, $0<u<1$, have strictly positive densities
on $(0,1)$ and thus have the same null sets. The support assertion follows.
\end{proof}

\begin{proposition}[Interval equilibrium and its energy]\label{prop:interval}
For $0<t<1$, $\beta_t$ maximizes $\mathcal E_t$ among all probability
measures on $[0,1]$, and
\begin{equation}\label{eq:interval-energy}
 e_t=U_t([0,1])
 =\frac{\Gamma((1+t)/2)\Gamma(1-t/2)}{2^t\sqrt\pi}.
\end{equation}
Its potential $\int_0^1|x-y|^t\dd\beta_t(y)$ equals $e_t$ at every
$x\in[0,1]$.
\end{proposition}

\begin{proof}
For any probability measure $\mu$ on $[0,1]$, Lemma~\ref{lem:quantile}
gives a monotone map $T$ taking values in $[0,1]$ almost everywhere.
Thus $L_T\leq1$, and~\eqref{eq:moment} and~\eqref{eq:mean-slope} give
\[
 \frac{\mathcal E_t(\mu)}{e_t}
 =\int D_T^t\dd P_t\leq L_T^t\leq1.
\]
This proves the equilibrium assertion without using a preexisting
interval-capacity formula.

Let $V_t(x)=\int|x-y|^t\dd\beta_t(y)$. For $x\in[0,1]$, maximality
applied to $(1-\varepsilon)\beta_t+\varepsilon\delta_x$ gives
\[
 0\geq\left.\frac{d}{d\varepsilon}
 \mathcal E_t\bigl((1-\varepsilon)\beta_t+\varepsilon\delta_x\bigr)
 \right|_{\varepsilon=0+}=2(V_t(x)-e_t).
\]
The function $V_t$ is continuous, $\beta_t$ has full support on $[0,1]$,
and $\int V_t\dd\beta_t=e_t$. Therefore $V_t(x)=e_t$ everywhere on
$[0,1]$. At zero this gives, with $a=(1-t)/2$,
\[
 e_t=\frac{B(a+t,a)}{B(a,a)}
 =\frac{\Gamma((1+t)/2)\Gamma(1-t)}{\Gamma((1-t)/2)}.
\]
The gamma duplication formula yields~\eqref{eq:interval-energy}.
\end{proof}

\begin{proof}[Proof of Theorem~\ref{thm:main}]
Fix $0<r<s<1$ and choose an $r$-energy maximizing measure $\mu$ on $K$.
By Lemma~\ref{lem:quantile}, choose $T$ with $T\push\beta_r=\mu$;
the measure $T\push\beta_s$ is also supported on $K$.
The transport comparison and Proposition~\ref{prop:interval} give
\begin{align*}
 \frac{C_r(K)}{C_r([0,1])}
 &=\left[\frac{\mathcal E_r(T\push\beta_r)}{e_r}\right]^{1/r}\\
 &\leq\left[\frac{\mathcal E_s(T\push\beta_s)}{e_s}\right]^{1/s}
 \leq\frac{C_s(K)}{C_s([0,1])}.
\end{align*}
The capacities in the denominators and $C_r(K)$ are positive, so this
is the desired ratio inequality. Translating a set does not change its
energy, and scaling it by $b>0$ multiplies $C_t$ by $b$. Thus $[0,1]$
and $[-1,1]$ have the same ratio; scaling~\eqref{eq:interval-energy}
also gives the gamma formula in~\eqref{eq:main}.

For equality, both steps of the displayed comparison must be equalities.
Theorem~\ref{thm:transport} gives $T(x)=A+Bx$ almost everywhere.
Here $B>0$, because $U_r(K)>0$. Thus $\mu$ is the affine image of
$\beta_r$, and its support is the interval $J=[A,A+B]\subset K$.
Its potential equals $U_r(J)=U_r(K)$ on $J$ by
Proposition~\ref{prop:interval} and scaling. If $K$ contained a point
$z$ outside $J$, choose the endpoint $b=A+B$ when $z>A+B$, and the
endpoint $b=A$ when $z<A$. For every $y\in J$, we have
$|z-y|>|b-y|$. Hence the potential at $z$ would strictly exceed the
potential at $b$, namely $U_r(K)$.
Adding a sufficiently small positive mass at $z$, as in the preceding
variation calculation, would increase the energy, a contradiction.
Consequently $K=J$. Conversely, every nondegenerate closed interval
attains equality by translation and scaling.
\end{proof}

\section{Examples, attribution, and verification}

For two points at distance $d>0$, weights $a$ and $1-a$ give energy
$2a(1-a)d^t$, maximized at $a=1/2$. Therefore
\[
 C_t(\{x,y\})=d\,2^{-1/t},\qquad
 \frac{C_s(\{x,y\})}{C_r(\{x,y\})}=2^{1/r-1/s}.
\]
Theorem~\ref{thm:main} makes this strictly greater than the interval
ratio. More generally the inequality is strict for every finite set
with at least two points, every disconnected compact set, and every
compact set that is not an interval. For an affine map $T(x)=A+Bx$ with
$B\geq0$,
the transport quantities satisfy $H_t(T)=B$, illustrating equality
in the stronger transport theorem.

The conjecture and the classical interval benchmark are attributed to
Clark and Laugesen~\cite{CL} and Fan~\cite{Fan}. The contribution here
is the monotone transport comparison, its Poisson-partition proof, and
its application to the full interval conjecture with equality
classification. The proof uses standard beta integrals, Poisson
superposition, likelihood-ratio ordering, and Jensen's inequality;
the required identities and comparison arguments are included above.
The accompanying literature review records the sources and public
repositories checked for earlier resolutions.

The proof is analytic and requires no numerical certificate. Optional
verification code checks the mixture normalization and finite
quantile-transport examples. The manuscript and its independent internal
mathematical checks were produced by GPT-6 Astra. The package contains
the LaTeX source and the proof audit; it makes no claim of proof-assistant
formalization or external peer review.

\begin{thebibliography}{9}
\bibitem{CL}
C. Clark and R. S. Laugesen,
\emph{Maximizing Riesz capacity ratios: conjectures and theorems},
SIAM Journal on Mathematical Analysis \textbf{57} (2025), 5683--5717.
\href{https://doi.org/10.1137/24M171992X}{doi:10.1137/24M171992X}.
Preprint: \href{https://arxiv.org/abs/2410.14809}{arXiv:2410.14809}.

\bibitem{Fan}
Q. Fan,
\emph{Riesz capacity ratios with negative exponents},
\href{https://arxiv.org/abs/2609.11186v1}{arXiv:2609.11186v1} (2026).
\end{thebibliography}
\end{document}
